\section{Algebraic extensions}
\label{section-algebraic-extensions}

\noindent
An important class of extensions are those where every element generates
a finite extension.

\begin{definition}
\label{definition-algebraic}
Consider a field extension $F/E$. An element $\alpha \in F$ is said to be
{\it algebraic} over $E$ if $\alpha$ is the root of some nonzero polynomial
with coefficients in $E$. If all elements of $F$ are algebraic then $F$ is
said to be an {\it algebraic extension} of $E$.
\end{definition}

\noindent
By Lemma \ref{lemma-field-extension-generated-by-one-element}, the
subextension $E(\alpha)$ is isomorphic either to the rational function
field $E(t)$ or to a quotient ring $E[t]/(P)$ for $P \in E[t]$ an
irreducible polynomial. In the latter case, $\alpha$ is algebraic over
$E$ (in fact, the proof of
Lemma \ref{lemma-field-extension-generated-by-one-element}
shows that we can pick $P$ such that $\alpha$ is a root of $P$);
in the former case, it is not.

\begin{example}
\label{example-C-algebraic-over-R}
The field $\mathbf{C}$ is algebraic over $\mathbf{R}$. Namely, if
$\alpha = a + ib$ in $\mathbf{C}$, then $\alpha^2 - 2a\alpha + a^2 + b^2 = 0$
is a polynomial equation for $\alpha$ over $\mathbf{R}$.
\end{example}

\begin{example}
\label{example-compact-riemann-surface-is-finite-over-P1}
Let $X$ be a compact Riemann surface, and let
$f \in \mathbf{C}(X) - \mathbf{C}$ any nonconstant meromorphic function
on $X$ (see Example \ref{example-field-of-meromorphic-functions}). Then it is
known that $\mathbf{C}(X)$ is algebraic over the subextension
$\mathbf{C}(f)$ generated by $f$. We shall not prove this.
\end{example}

\begin{lemma}
\label{lemma-algebraic-goes-up}
Let $K/E/F$ be a tower of field extensions.
\begin{enumerate}
\item If $\alpha \in K$ is algebraic over $F$, then $\alpha$ is algebraic
over $E$.
\item If $K$ is algebraic over $F$, then $K$ is algebraic over $E$.
\end{enumerate}
\end{lemma}

\begin{proof}
This is immediate from the definitions.
\end{proof}

\noindent
We now show that there is a deep connection between finiteness and being
algebraic.

\begin{lemma}
\label{lemma-finite-is-algebraic}
A finite extension is algebraic. In fact, an extension $E/k$ is algebraic
if and only if every subextension $k(\alpha)/k$ generated by some
$\alpha \in E$ is finite.
\end{lemma}

\noindent
In general, it is very false that an algebraic extension is finite.

\begin{proof}
Let $E/k$ be finite, say of degree $n$. Choose $\alpha \in E$. Then the
elements $1, \alpha, \ldots, \alpha^n$ are linearly
dependent over $k$, or we would necessarily have $[E : k] > n$. A relation of
linear dependence now gives the desired polynomial that $\alpha$ must satisfy.

\medskip\noindent
For the last assertion, note that a monogenic extension $k(\alpha)/k$ is
finite if and only if $\alpha$ is algebraic over $k$, by
Examples \ref{example-degree-rational-function-field} and
\ref{example-degree-simple-algebraic-extension}.
So if $E/k$ is algebraic, then each $k(\alpha)/k$, $\alpha \in E$, is a finite
extension, and conversely.
\end{proof}

\noindent
We can extract a lemma of the last proof (really of
Examples \ref{example-degree-rational-function-field} and
\ref{example-degree-simple-algebraic-extension}):
a monogenic extension is finite if and only if it is algebraic.
We shall use this observation in the next result.

\begin{lemma}
\label{lemma-algebraic-finitely-generated}
\begin{slogan}
A finitely generated algebraic extension is finite.
\end{slogan}
Let $k$ be a field, and let $\alpha_1, \alpha_2, \ldots, \alpha_n$ be elements
of some extension field such that each $\alpha_i$ is algebraic over $k$. Then
the extension $k(\alpha_1, \ldots, \alpha_n)/k$ is finite.
That is, a finitely generated algebraic extension is finite.
\end{lemma}

\begin{proof}
Indeed, each extension
$k(\alpha_{1}, \ldots, \alpha_{i+1})/k(\alpha_1, \ldots, \alpha_{i})$
is generated by one element and algebraic, hence finite.
By multiplicativity of degree (Lemma \ref{lemma-multiplicativity-degrees})
we obtain the result.
\end{proof}

\noindent
The set of complex numbers that are algebraic over $\mathbf{Q}$ are simply
called the {\it algebraic numbers.} For instance, $\sqrt{2}$ is algebraic,
$i$ is algebraic, but $\pi$ is not.
It is a basic fact that the algebraic numbers form a field, although it is not
obvious how to prove this from the definition that a number is algebraic
precisely when it satisfies a nonzero polynomial equation with rational
coefficients (e.g. by polynomial equations).

\begin{lemma}
\label{lemma-algebraic-elements}
Let $E/k$ be a field extension. Then the elements of $E$ algebraic over $k$
form a subextension of $E/k$.
\end{lemma}

\begin{proof}
Let $\alpha, \beta \in E$ be algebraic over $k$. Then $k(\alpha, \beta)/k$
is a finite extension by Lemma \ref{lemma-algebraic-finitely-generated}.
It follows that $k(\alpha + \beta) \subset k(\alpha, \beta)$ is a finite
extension, which implies that $\alpha + \beta$ is algebraic by
Lemma \ref{lemma-finite-is-algebraic}. Similarly for the difference,
product and quotient of $\alpha$ and $\beta$.
\end{proof}

\noindent
Many nice properties of field extensions, like those of rings, will have the
property that they will be preserved by towers and composita.

\begin{lemma}
\label{lemma-algebraic-permanence}
Let $E/k$ and $F/E$ be algebraic extensions of fields. Then $F/k$ is an
algebraic extension of fields.
\end{lemma}

\begin{proof}
Choose $\alpha \in F$. Then $\alpha$ is algebraic over $E$.
The key observation is that $\alpha$ is algebraic over a
finitely generated subextension of $k$.
That is, there is a finite set $S \subset E$ such that $\alpha $ is algebraic
over $k(S)$: this is clear because being algebraic means that a certain
polynomial in $E[x]$ that $\alpha$ satisfies exists, and as $S$ we can take the
coefficients of this polynomial. It follows that $\alpha$ is algebraic over
$k(S)$. In particular, the extension $k(S, \alpha)/ k(S)$ is finite.
Since $S$ is a finite set, and $k(S)/k$ is algebraic,
Lemma \ref{lemma-algebraic-finitely-generated} shows that
$k(S)/k$ is finite. Using multiplicativity
(Lemma \ref{lemma-multiplicativity-degrees})
we find that $k(S,\alpha)/k$ is finite, so $\alpha$ is algebraic over $k$.
\end{proof}

\noindent
The method of proof in the previous argument --- that being algebraic
over $E$ was a property that {\it descended} to a finitely generated
subextension of $E$ --- is an idea that recurs throughout algebra.
It often allows one to reduce general commutative algebra questions
to the Noetherian case for example.

\begin{lemma}
\label{lemma-size-algebraic-extension}
Let $E/F$ be an algebraic extension of fields. Then the cardinality $|E|$
of $E$ is at most $\max(\aleph_0, |F|)$.
\end{lemma}

\begin{proof}
Let $S$ be the set of nonconstant polynomials with coefficients in $F$.
For every $P \in S$ the set of roots
$r(P, E) = \{\alpha \in E \mid P(\alpha) = 0\}$
is finite (details omitted). Moreover, the fact that $E$ is algebraic
over $F$ implies that $E = \bigcup_{P \in S} r(P, E)$.
It is clear that $S$ has cardinality bounded by $\max(\aleph_0, |F|)$
because it is a countable union of finite products of copies of $F$.
Thus so does $E$.
\end{proof}

\begin{lemma}
\label{lemma-subalgebra-algebraic-extension-field}
Let $E/F$ be a finite or more generally an algebraic extension of fields.
Any subring $F \subset R \subset E$ is a field.
\end{lemma}

\begin{proof}
Let $\alpha \in R$ be nonzero. Then $1, \alpha, \alpha^2, \ldots$
are contained in $R$. By Lemma \ref{lemma-finite-is-algebraic}
we find a nontrivial relation
$a_0 + a_1 \alpha + \ldots + a_d \alpha^d = 0$ with $a_i \in F$.
We may assume $a_0 \not = 0$ because if not we can divide the relation
by $\alpha$ to decrease $d$. Then we see that
$$
a_0 = \alpha (- a_1  - \ldots - a_d \alpha^{d - 1})
$$
which proves that the inverse of $\alpha$ is the element
$a_0^{-1} (- a_1  - \ldots - a_d \alpha^{d - 1})$
of $R$.
\end{proof}

\begin{lemma}
\label{lemma-algebraic-extension-self-map}
Let $E/F$ an algebraic extension of fields. Any $F$-algebra map
$f : E \to E$ is an automorphism.
\end{lemma}

\begin{proof}
If $E/F$ is finite, then $f : E \to E$ is an $F$-linear
injective map (Lemma \ref{lemma-field-maps-injective})
of finite dimensional vector spaces, and hence bijective.
In general we still see that $f$ is injective.
Let $\alpha \in E$ and let $P \in F[x]$ be a
polynomial such that $P(\alpha) = 0$.
Let $E' \subset E$ be the subfield of $E$ generated
by the roots $\alpha = \alpha_1, \ldots, \alpha_n$ of $P$ in $E$.
Then $E'$ is finite over $F$ by Lemma \ref{lemma-algebraic-finitely-generated}.
Since $f$ preserves the set of roots, we find that
$f|_{E'} : E' \to E'$. Hence $f|_{E'}$ is an isomorphism
by the first part of the proof and we conclude that $\alpha$
is in the image of $f$.
\end{proof}






